\chapter[METHODOLOGY]{Methodology}
\label{ch:methodology}

\section{Research Design}
\label{sec:research-design}
% -----------------------------------------------------------------
% Overall approach and why it fits the questions. A short paragraph
% plus an overall workflow figure is usually enough here.
% -----------------------------------------------------------------

\note{Research design to be drafted}

% \begin{figure}[htbp]
%   \centering
%   \includegraphics[width=\textwidth]{workflow.pdf}
%   \caption{Overall research workflow.}
%   \label{fig:workflow}
% \end{figure}


\section{Project Activities}
\label{sec:project-activities}
% -----------------------------------------------------------------
% Walk through the pipeline, one subsection per stage. Each stage
% should name the tools and state what "done" looks like.
% -----------------------------------------------------------------

\subsection{Data Acquisition}
% Sensor selection, rig, sample surfaces, controlled vs field data.
% Say honestly whether data is collected or sourced -- the datasets
% section of Chapter 2 makes this the biggest risk in the project.
\note{Data acquisition to be drafted}

\subsection{Calibration and Registration}
% Targetless extrinsic calibration; a printed chessboard is
% invisible to thermal, so state the calibration target used.
\note{Calibration and registration to be drafted}

\subsection{Preprocessing and Annotation}
\note{Preprocessing and annotation to be drafted}

\subsection{Model Development}
% Baseline (RGB only), then the fused model.
\note{Model development to be drafted}

\subsection{Ablation Protocol}
% RGB -> RGB + thermal -> full trimodal, with the metrics fixed in
% advance. This is the methodological contribution; give it space
% and state the decision rule for "the modality earned its place"
% BEFORE any results exist.
\note{Ablation protocol to be drafted}

\subsection{Validation and Benchmarking}
\note{Validation and benchmarking to be drafted}


\section{Objectives and Key Milestones}
\label{sec:milestones}
% -----------------------------------------------------------------
% The Word template carries a SmartArt diagram mapping three
% objectives onto five milestones. The TikZ figure below reproduces
% that mapping -- edit the labels once the objectives are settled.
% -----------------------------------------------------------------

\begin{figure}[htbp]
\centering
\begin{tikzpicture}[
  obj/.style={draw,thick,rounded corners,align=center,
              minimum width=3.4cm,minimum height=1.15cm,
              fill=blue!8,inner sep=3pt},
  ms/.style={draw,thick,rounded corners,align=center,
             minimum width=5.0cm,minimum height=0.9cm,
             fill=orange!12,inner sep=3pt},
  arr/.style={-{Latex[length=2mm]},thick,gray!65}
]
% --- milestones: an evenly spaced ladder on the right ---
\node[ms] (m1) at (0,0) {\textbf{Milestone 1} \footnotesize \note{label}};
\node[ms,below=0.32cm of m1] (m2) {\textbf{Milestone 2} \footnotesize \note{label}};
\node[ms,below=0.32cm of m2] (m3) {\textbf{Milestone 3} \footnotesize \note{label}};
\node[ms,below=0.32cm of m3] (m4) {\textbf{Milestone 4} \footnotesize \note{label}};
\node[ms,below=0.32cm of m4] (m5) {\textbf{Milestone 5} \footnotesize \note{label}};

% --- objectives: centred against the milestones each one owns ---
\node[obj] (o1) at ($(m1.west)!0.5!(m2.west) + (-2.5cm,0)$)
  {\textbf{Objective 1}\\[-2pt]\footnotesize \note{short label}};
\node[obj] (o2) at ($(m3.west) + (-2.5cm,0)$)
  {\textbf{Objective 2}\\[-2pt]\footnotesize \note{short label}};
\node[obj] (o3) at ($(m4.west)!0.5!(m5.west) + (-2.5cm,0)$)
  {\textbf{Objective 3}\\[-2pt]\footnotesize \note{short label}};

\draw[arr] (o1.east) -- (m1.west);
\draw[arr] (o1.east) -- (m2.west);
\draw[arr] (o2.east) -- (m3.west);
\draw[arr] (o3.east) -- (m4.west);
\draw[arr] (o3.east) -- (m5.west);
\end{tikzpicture}
\caption{Mapping of research objectives to key milestones.}
\label{fig:objectives-milestones}
\end{figure}

\begin{table}[htbp]
\centering
\caption{Key milestones and expected completion.}
\label{tab:milestones}
\begin{tabularx}{\textwidth}{@{}l X l@{}}
\toprule
\textbf{Milestone} & \textbf{Description} & \textbf{Target} \\
\midrule
M1 & \note{description} & \note{date} \\
M2 & \note{description} & \note{date} \\
M3 & \note{description} & \note{date} \\
M4 & \note{description} & \note{date} \\
M5 & \note{description} & \note{date} \\
\bottomrule
\end{tabularx}
\end{table}


\section{Study Plan}
\label{sec:study-plan}
% -----------------------------------------------------------------
% The Gantt chart, May 2026 - April 2028. Either include the
% existing PDF as a figure, or redraw it with pgfgantt for a
% cleaner result. Note that the existing chart still carries the
% superseded pitting-specific title -- fix that before it goes in.
% -----------------------------------------------------------------

\note{Gantt chart to be inserted}

% \begin{figure}[htbp]
%   \centering
%   \includegraphics[width=\textwidth]{gantt.pdf}
%   \caption{Study plan, May 2026 to April 2028.}
%   \label{fig:gantt}
% \end{figure}


\section{Preliminary Results}
\label{sec:preliminary-results}
% -----------------------------------------------------------------
% The template names this section explicitly (3.1 in the original),
% so it cannot be dropped. Even a small proof of concept -- a
% baseline RGB detector on a public dataset, or a registration
% trial on two modalities -- is worth far more here than nothing.
% -----------------------------------------------------------------

\note{Preliminary results to be drafted}
